\documentclass[11pt]{amsart}
\usepackage[margin=1.15in]{geometry}
\usepackage{amsmath,amssymb,amsthm}
\usepackage[hidelinks]{hyperref}

\newtheorem{theorem}{Theorem}
\newtheorem{lemma}[theorem]{Lemma}
\newtheorem{proposition}[theorem]{Proposition}
\newtheorem{corollary}[theorem]{Corollary}
\theoremstyle{remark}
\newtheorem{remark}[theorem]{Remark}

\newcommand{\E}{\mathbb{E}}
\renewcommand{\P}{\mathbb{P}}
\newcommand{\R}{\mathbb{R}}
\newcommand{\cF}{\mathcal{F}}

\title[Loop-erased ants need not find shortest paths]{Loop-erased ants need not find shortest paths}
\author{Yitzchak Shmalo}
\address{Einstein Institute of Mathematics, The Hebrew University of Jerusalem,\newline Givat Ram, Jerusalem, Israel}

\begin{document}
\begin{abstract}
In the ant process on a finite graph, each ant performs a random walk from a nest to a food source with transition probabilities proportional to the current edge weights, and every edge on the loop erasure of its trajectory gains weight one. Kious, Mailler and Schapira conjectured that the normalised weight of every edge outside the shortest paths tends to zero when the trajectory is reversed before the loops are erased. Armfield and Mailler conjectured the same for the erasure of the forward trajectory and proved it when the nest and the food are at distance two. By the Aldous--Broder and Wilson algorithms the two erasure rules give the reinforced path the same law, so the two processes coincide in law. We then give a graph with six vertices on which, with positive probability, an edge on no shortest path keeps a limiting normalised weight in $[2499/20000,1/8]$. Adding a path at the nest gives counterexamples at every distance at least three.
\end{abstract}

\subjclass[2020]{Primary 60K35; Secondary 60J10, 62L20, 05C81}
\keywords{reinforced random walk, ant process, loop-erased random walk, uniform spanning tree, stochastic approximation}

\maketitle

\section{Introduction}

Let $G$ be a finite connected simple graph with edge set $E$ and two distinct marked vertices, the nest $N$ and the food $F$. The ant process is a Markov chain $(W(n))_{n\ge0}$ of positive edge weights with $W_e(0)=1$ for every $e\in E$. Given $W(n)$, the $n$-th ant performs the random walk $\omega^{(n)}$ on the network $(G,W(n))$, started at $N$ and stopped at its first visit to $F$. From a vertex $x$ it moves to a neighbour $y$ with probability $W_{xy}(n)/\sum_{y'}W_{xy'}(n)$, the sum running over the neighbours $y'$ of $x$. Let $\Gamma_n$ be the self-avoiding path from $N$ to $F$ that one of the two rules below extracts from $\omega^{(n)}$. Then
\[
W_e(n+1)=W_e(n)+\mathbf 1\{e\in\Gamma_n\},\qquad e\in E,
\]
where $\mathbf 1\{e\in\Gamma_n\}$ is $1$ if $e$ lies on $\Gamma_n$ and $0$ otherwise.
Two rules for $\Gamma_n$ are studied in the literature.
\begin{itemize}
\item The backward rule, used in the loop-erased ant process of Kious, Mailler and Schapira \cite{KMS22a}: $\Gamma_n$ is obtained by following the trajectory backwards from $F$ and erasing loops. It is the chronological loop erasure of the time reversal of $\omega^{(n)}$, read from $N$ to $F$.
\item The forward rule, used in the forward-loop-erased ant process of Armfield and Mailler \cite{AM26}: $\Gamma_n$ is the chronological loop erasure of $\omega^{(n)}$ itself.
\end{itemize}
We call $W_e(n)/n$ the normalised weight of $e$, and we say that the ants find the shortest paths if $W_e(n)/n\to0$ almost surely for every edge $e$ that lies on no shortest path from $N$ to $F$.

Kious, Mailler and Schapira conjectured \cite[Conjecture~1.1]{KMS22a} that in the loop-erased process $W(n)/n$ converges almost surely and that the ants find the shortest paths. They proved this for series-parallel graphs with source $N$ and sink $F$ \cite[Theorem~1.3]{KMS22a}. Armfield and Mailler proved that in the forward-loop-erased process the ants find the shortest paths whenever $N$ and $F$ are at distance two, with no other assumption on $G$ \cite[Theorem~1.1]{AM26}, and conjectured that this holds at every distance \cite[Section~1]{AM26}. Their proof uses the fact that, by Wilson's algorithm \cite{Wil96}, the path reinforced under the forward rule is the path from $N$ to $F$ in a weighted spanning tree, a description they regard as particular to the forward rule. The Aldous--Broder algorithm \cite{Ald90,Bro89} gives the same description for the backward rule, so the two processes have the same law (Proposition~\ref{prop:same}), and \cite[Theorem~1.1]{AM26} holds for the loop-erased process as well.

At distance three both conjectures fail. Let $G_6$ consist of the paths $0\,1\,2\,5$ and $0\,3\,4\,5$ from the nest $0$ to the food $5$, joined by the edge $13$, and write $u$, $v$, $z$ for $W_{01}$, $W_{12}$ and $W_{13}$ divided by $n+2$. Let $f=(f_u,f_v,f_z)$ be the mean drift of $(u,v,z)$: $f_u$ is the probability that the next ant reinforces $01$, minus $u$, and similarly for $f_v$ and $f_z$. The drift vanishes on the whole curve
\[
u=v=x,\qquad z=\frac{x(1-x)}2,\qquad 0<x<1,
\]
along which the edge $13$, which lies on no shortest path, is reinforced exactly in proportion to its own weight. The linearisation at $(1/2,1/2,1/8)$ has a zero eigenvalue along the curve, so the stability criteria of stochastic approximation, which ask for a linearly stable equilibrium, do not apply. Put $q=(u+2v)/3$, $\delta=u-v$, $r=z-q(1-q)/2$ and $V=\delta^2+r^2$, so that the curve is $\{\delta=r=0\}$ and $q$ is the position along it. Then
\[
\nabla V\cdot f\le-\tfrac18\,V,\qquad |\nabla q\cdot f|\le V\qquad\text{when } |q-\tfrac12|,\ |\delta|,\ |r|\le\tfrac1{100}\ \text{and}\ V\le10^{-4}.
\]
The first inequality draws the process to the curve at a linear rate. By the second, the position $q$ moves only by an amount of order $V$ plus a martingale. A process that comes close enough to $(1/2,1/2,1/8)$ at a late enough time therefore stays near the curve forever with positive probability and converges to a random point $(X,X,X(1-X)/2)$ of it with $|X-\tfrac12|\le\tfrac1{100}$, and then $W_{13}(n)/n$ tends to $X(1-X)/2\ge2499/20000$ (Theorem~\ref{thm:main}).

The simplest graph that is not series-parallel is the losange, with vertices $N,a,b,F$ and edges $Na$, $Nb$, $aF$, $bF$, $ab$. On it Kious, Mailler and Schapira proved the corresponding statement for the uniform-geodesic model, in which each ant reinforces a shortest path from $N$ to $F$, chosen uniformly, in the subgraph it has traversed \cite[Theorem~1.4]{KMS22a}; for the loop-erased process on the losange, Proposition~\ref{prop:same} and \cite[Theorem~1.1]{AM26} show that the ants find the shortest paths. When $N$ and $F$ are adjacent, the argument of \cite[Proposition~2.8]{KMS22b}, written for the trace-reinforced process, shows that the ants find the shortest paths in every ant process whose reinforced path is self-avoiding, as noted in \cite[Section~1]{AM26}; this covers both rules. For the loop-erased process with several nests, see Mailler and Varin \cite{MV26}.

Both spanning-tree algorithms are classical, and together they give the following.

\begin{proposition}\label{prop:same}
For all positive weights on $G$, the reinforced path $\Gamma$ has the same law under the two rules. In both cases it is the path from $N$ to $F$ in a spanning tree of $G$ chosen with probability proportional to the product of its edge weights. Consequently the loop-erased and the forward-loop-erased ant processes on $(G,N,F)$ have the same law.
\end{proposition}

The equality in law of the two reinforced paths is a form of the invariance of loop-erased random walk under time reversal, which goes back to Lawler; see \cite[Corollary~4.4 and the remark before it]{LP16}.

In full, $G_6$ has vertex set $\{0,1,2,3,4,5\}$, nest $0$, food $5$ and edges
\[
01,\quad 12,\quad 25,\quad 03,\quad 34,\quad 45,\quad 13 .
\]
The paths $0\,1\,2\,5$ and $0\,3\,4\,5$, of length three, are the shortest paths from $0$ to $5$. We call $13$ the cross edge and the other six edges the branch edges. The two other self-avoiding paths from $0$ to $5$, namely $0\,1\,3\,4\,5$ and $0\,3\,1\,2\,5$, have length four and are the only ones that use $13$; we call them the crossing paths.

\begin{theorem}\label{thm:main}
For the ant process on $G_6$, under either rule, with positive probability
\[
\frac1n\bigl(W_{01},W_{12},W_{25},W_{03},W_{34},W_{45},W_{13}\bigr)(n)\to\Bigl(X,X,X,1-X,1-X,1-X,\frac{X(1-X)}2\Bigr)
\]
for a random variable $X$ with values in $[49/100,51/100]$.
On this event $\lim_n W_{13}(n)/n\in[2499/20000,\,1/8]$, although $13$ lies on no shortest path. Hence the ants do not find the shortest paths on $G_6$.
\end{theorem}

\begin{corollary}\label{cor:distance}
For every integer $d\ge3$ there is a finite graph on which $N$ and $F$ are at distance $d$ and, under either rule, with positive probability an edge that lies on no shortest path has a positive limiting normalised weight.
\end{corollary}

A heuristic explains the mechanism at the point where the branch weights are $4$ and the cross weight is $1$. The two crossing paths are chosen with probability $1/16$ each and traverse $13$ in opposite directions. In the electrical network with unit voltage at $0$ and zero voltage at $5$, the vertices $1$ and $3$ are both at voltage $2/3$, so no current flows through $13$. When a unit current flows from $0$ to $5$, the expected number of signed crossings of an edge by the walk is the current through it \cite[Proposition~2.2]{LP16}, so the expected numbers of crossings of $13$ in the two directions are equal. The reinforcement counts both directions, and by \eqref{eq:paths} at $u=v$ the probability that the ant reinforces $13$ is $1/8$. This is exactly $W_{13}/(n+2)$ at this point, so $13$ is reinforced in proportion to its own weight and neither grows nor decays relative to the others.

Section~\ref{sec:same} proves Proposition~\ref{prop:same}. Section~\ref{sec:drift} computes the expected reinforcement on $G_6$, Section~\ref{sec:trap} proves Theorem~\ref{thm:main}, and Section~\ref{sec:cor} proves Corollary~\ref{cor:distance} and shows that six vertices are needed (Remark~\ref{rem:minimal}).

\section{The two erasure rules}\label{sec:same}

Fix positive weights $w$ on $G$. For a path $\omega=(\omega_0,\dots,\omega_m)$ let $\mathrm{LE}(\omega)$ be its chronological loop erasure: $\sigma_0=\omega_0$ and, as long as $\sigma_i\ne\omega_m$, $\sigma_{i+1}=\omega_{j+1}$, where $j$ is the last time at which $\omega$ visits $\sigma_i$. For a spanning tree $T$ let $\pi_{x,y}(T)$ be the path from $x$ to $y$ in $T$, and let $\mathsf T$ be a spanning tree chosen with probability proportional to $\prod_{e\in T}w_e$.

By Wilson's algorithm \cite{Wil96} (see \cite[Corollary~4.3]{LP16}), for any vertices $x\ne y$ the loop erasure of the random walk from $x$ stopped at its first visit to $y$ has the law of $\pi_{x,y}(\mathsf T)$. With $x=N$ and $y=F$ this is the statement of Proposition~\ref{prop:same} for the forward rule.

For the backward rule, let $(Z_k)_{k\ge0}$ be the random walk on $(G,w)$ started at $N$, let $m$ be its first visit to $F$, write $\omega=(Z_0,\dots,Z_m)$ and let $\omega^R=(\omega_m,\dots,\omega_0)$ be its reversal. The last visit of $\omega^R$ to a vertex is the first visit of $\omega$ to it. Hence, writing $\mathrm{LE}(\omega^R)=(\sigma_0,\sigma_1,\dots)$, we have $\sigma_0=F$ and $\sigma_{i+1}$ is the vertex from which $\omega$ first enters $\sigma_i$. So the backward rule follows, from $F$, the edges by which $\omega$ first entered each vertex.

For $x\ne N$ let $H(x)$ be the first time the walk $(Z_k)_{k\ge0}$ visits $x$; it is finite almost surely. The edges $\{x,Z_{H(x)-1}\}$, $x\ne N$, form a spanning tree $T_{\mathrm{AB}}$, and by the Aldous--Broder algorithm \cite{Ald90,Bro89} (see \cite[Corollary~4.9]{LP16}) $T_{\mathrm{AB}}$ has the law of $\mathsf T$. The vertices on the path followed by the backward rule are visited by time $m$, and their first-entrance edges are edges of $T_{\mathrm{AB}}$, so the path followed by the backward rule is the path from $F$ to $N$ in $T_{\mathrm{AB}}$. Read from $N$ to $F$, it is $\pi_{N,F}(T_{\mathrm{AB}})$, which has the law of $\pi_{N,F}(\mathsf T)$.

Since the transition law of $W(n)$ depends on the rule only through the law of $\Gamma_n$ given $W(n)$, the two processes have the same law. \qed

\section{The expected reinforcement on \texorpdfstring{$G_6$}{G6}}\label{sec:drift}

By Proposition~\ref{prop:same} we may work with either rule. Every reinforced path is one of
\[
U=0\,1\,2\,5,\qquad L=0\,3\,4\,5,\qquad C=0\,1\,3\,4\,5,\qquad B=0\,3\,1\,2\,5 .
\]
Each of them uses exactly one of $01,03$ and exactly one of $12,34$; the edges $12$ and $25$ lie on the same paths, $U$ and $B$, and so do $34$ and $45$, on $L$ and $C$. Hence, for every $n$,
\[
W_{12}=W_{25},\qquad W_{34}=W_{45},\qquad W_{01}+W_{03}=W_{12}+W_{34}=n+2 .
\]
Put
\[
u=\frac{W_{01}}{n+2},\qquad v=\frac{W_{12}}{n+2},\qquad z=\frac{W_{13}}{n+2}.
\]
The other weights divided by $n+2$ are $1-u$ for $03$, $v$ for $25$, and $1-v$ for $34$ and $45$. The walk does not change when all weights are multiplied by the same constant, so the law of $\Gamma_n$ depends only on $Y_n=(u_n,v_n,z_n)\in(0,1)^3$.

The graph $G_6$ has fourteen spanning trees. A tree whose path from $0$ to $5$ is $U$ contains $01,12,25$ and joins $3$ and $4$ by one of the pairs $\{03,34\},\{03,45\},\{34,45\},\{13,34\},\{13,45\}$; a tree whose path is $C$ contains $01,13,34,45$ and attaches $2$ by $12$ or by $25$; and symmetrically for $L$ and $B$. After division by the common factor $v(1-v)$ of their weights, Proposition~\ref{prop:same} gives
\begin{equation}\label{eq:paths}
\begin{aligned}
p_U&=\frac{uv\,(2(1-u)+(1-v)+2z)}{D},&\qquad p_L&=\frac{(1-u)(1-v)\,(2u+v+2z)}{D},\\
p_C&=\frac{2uz(1-v)}{D},&\qquad p_B&=\frac{2(1-u)zv}{D},
\end{aligned}
\end{equation}
where
\[
D=2u(1-u)+v(1-v)+2z
\]
is the sum of the four numerators. At the point $u=v=1/2$, $z=1/8$ the four probabilities are $7/16$, $7/16$, $1/16$ and $1/16$.

The edge $01$ lies on $U$ and $C$, the edge $12$ on $U$ and $B$, and $13$ on $C$ and $B$. Since $u_{n+1}=(W_{01}(n)+\mathbf 1\{01\in\Gamma_n\})/(n+3)$, and similarly for $v$ and $z$,
\begin{equation}\label{eq:sa}
Y_{n+1}-Y_n=\gamma_n\bigl(f(Y_n)+\xi_{n+1}\bigr),\qquad \gamma_n=\frac1{n+3},
\end{equation}
where $\cF_n$ is the $\sigma$-field generated by $W(0),\dots,W(n)$, the noise $\xi_{n+1}=(\xi^{(u)}_{n+1},\xi^{(v)}_{n+1},\xi^{(z)}_{n+1})$ satisfies $\E[\xi_{n+1}\mid\cF_n]=0$, every coordinate of $\xi_{n+1}$ has absolute value at most one, and $f=(f_u,f_v,f_z)=(p_U+p_C-u,\;p_U+p_B-v,\;p_C+p_B-z)$. From \eqref{eq:paths},
\begin{equation}\label{eq:drift}
f_u=-\frac{2u(1-u)(u-v)}{D},\qquad f_v=\frac{v(1-v)(u-v)}{D},\qquad f_z=\frac{z\,(2u^2-4uv+v^2+v-2z)}{D}.
\end{equation}
The zeros of $f$ in the interior include the curve $u=v=x$, $z=x(1-x)/2$, $0<x<1$. The Jacobian of $f$ at $(1/2,1/2,1/8)$ has eigenvalues $0,-3/4,-1/4$, the zero eigenvalue along the curve, so the linearisation alone does not decide whether the process can settle near the curve. We prove that it does with positive probability.

\section{Proof of Theorem~\ref{thm:main}}\label{sec:trap}

\subsection{Coordinates and a Lyapunov function}
Put
\[
q=\frac{u+2v}3,\qquad \delta=u-v,\qquad r=z-\frac{q(1-q)}2,\qquad V=\delta^2+r^2,
\]
so that $u=q+2\delta/3$, $v=q-\delta/3$ and the curve of zeros is $\{\delta=r=0\}$. For a differentiable function $h$ of $Y$ write $\mathcal L h=\nabla h\cdot f$. From \eqref{eq:drift},
\begin{equation}\label{eq:Lq}
\mathcal Lq=\frac{2\delta^2(u+v-1)}{3D},\qquad \mathcal L\delta=-\frac{2u(1-u)+v(1-v)}{D}\,\delta .
\end{equation}
Let $\eta=1/100$ and $t=q-1/2$, and let $\mathcal B$ be the set of $Y$ with $|t|\le\eta$, $|\delta|\le\eta$ and $|r|\le\eta$. Expanding $D\,\mathcal LV$ as a polynomial in $t,\delta,r$ gives
\begin{equation}\label{eq:DLV}
\begin{aligned}
D\,\mathcal LV={}&-\tfrac32\delta^2-\tfrac12r^2+2\delta^4+\tfrac49\delta^3rt+4\delta^3t+\tfrac{34}9\delta^2r^2+\tfrac79\delta^2rt^2+\tfrac{17}{36}\delta^2r+6\delta^2t^2\\
&+\tfrac43\delta r^2t-\tfrac23\delta rt^3+\tfrac16\delta rt-4r^3+2r^2t^2 .
\end{aligned}
\end{equation}
Every monomial $\lambda\,t^i\delta^jr^k$ after the first two has $j+k\ge2$, and on $\mathcal B$ with $V\le\eta^2$ its absolute value is at most $|\lambda|\,\eta^{i+j+k-2}V$. The sum of these bounds is $\frac{434717}{9000000}V<\frac14V$, so $D\,\mathcal LV\le-\frac14V$ there. On $\mathcal B$ we also have $\frac12\le D\le2$ and $|u+v-1|=|2t+\delta/3|\le\frac73\eta$. Therefore
\begin{equation}\label{eq:box}
\mathcal LV\le-\tfrac18V,\qquad |\mathcal Lq|\le\tfrac{7}{225}\,\delta^2\le V\qquad\text{on }\mathcal B\cap\{V\le\eta^2\}.
\end{equation}

\subsection{One step}
On $[0,1]^3$ we have $|r|\le1$, $\|\nabla\delta\|^2=2$, $\|\nabla q\|^2=5/9$, $\|\nabla r\|^2\le41/36$ and $\nabla^2 r=\nabla q\,\nabla q^{\top}$. Hence the operator norm of $\nabla^2V=2\nabla\delta\nabla\delta^\top+2\nabla r\nabla r^\top+2r\nabla^2r$ is at most $4+41/18+10/9=133/18<8$. By \eqref{eq:sa} a step has Euclidean length at most $\sqrt3\,\gamma_n$ and the segment it spans lies in $[0,1]^3$, so Taylor's theorem and \eqref{eq:box} give, whenever $Y_n\in\mathcal B$ and $V(Y_n)\le\eta^2$,
\begin{equation}\label{eq:step}
\E\bigl[V(Y_{n+1})-V(Y_n)\mid\cF_n\bigr]\le\gamma_n\,\mathcal LV(Y_n)+12\gamma_n^2\le-\tfrac18\gamma_nV(Y_n)+12\gamma_n^2 .
\end{equation}

\subsection{Trapping from a restart state}
Let $M\ge1$ be an integer and suppose that at some time all six branch weights equal $4M$ and $W_{13}=M$; we call this a restart state. Then $n+2=8M$ and $Y=(1/2,1/2,1/8)$, so $t=\delta=r=0$. Count time from this moment, write $Y_k$ for the state $k$ ants later and $S=8M$; the step sizes are $\gamma_k=1/(S+k+1)$, and $\sum_{k\ge0}\gamma_k^2\le1/S$.

Let $\tau$ be the first $k$ with $|q_k-\frac12|\ge\eta$ or $V(Y_k)\ge\eta^2$. Before $\tau$ the state lies in $\mathcal B$ with $V<\eta^2$, so \eqref{eq:step} applies to every step taken before $\tau$, including the step that reaches $\tau$. Let $K\ge0$ be an integer. Summing the increments stopped at $K\wedge\tau$,
\begin{equation}\label{eq:sum}
\E V(Y_{K\wedge\tau})+\frac18\,\E\sum_{k<K\wedge\tau}\gamma_kV(Y_k)\le\frac{12}S .
\end{equation}
Letting $K\to\infty$, Fatou's lemma and Markov's inequality give
\begin{equation}\label{eq:exitV}
\P\bigl(\tau<\infty,\ V(Y_\tau)\ge\eta^2\bigr)\le\frac{12}{S\eta^2},\qquad \E\sum_{k<\tau}\gamma_kV(Y_k)\le\frac{96}S .
\end{equation}
Since $q$ is linear,
\[
q_{K\wedge\tau}-\tfrac12=\sum_{k<K\wedge\tau}\gamma_k\,\mathcal Lq(Y_k)+\mathcal M_K,\qquad \mathcal M_K=\sum_{k<K\wedge\tau}\gamma_k\,\frac{\xi^{(u)}_{k+1}+2\xi^{(v)}_{k+1}}3 ,
\]
and $(\mathcal M_K)$ is a martingale whose increments have conditional variance at most $\gamma_k^2$. By \eqref{eq:box}, \eqref{eq:exitV} and Markov's inequality, $\P\bigl(\sum_{k<\tau}\gamma_k|\mathcal Lq(Y_k)|\ge\eta/2\bigr)\le192/(S\eta)$, and by Doob's inequality $\P(\sup_K|\mathcal M_K|\ge\eta/2)\le4/(S\eta^2)$. If $\tau<\infty$ and $V(Y_\tau)<\eta^2$, then $|q_\tau-\frac12|\ge\eta$ and one of these two events occurs. Hence
\begin{equation}\label{eq:trap}
\P(\tau=\infty)\ge1-\frac{16}{S\eta^2}-\frac{192}{S\eta}.
\end{equation}
For $M=50001$, so $S=400008$, the right-hand side equals $3943/7143>1/2$.

\subsection{\texorpdfstring{Convergence on $\{\tau=\infty\}$}{Convergence when the process stays trapped}}
By \eqref{eq:step}, $V(Y_{K\wedge\tau})+12\sum_{k\ge K}\gamma_k^2$ is a nonnegative supermartingale, so $V(Y_{K\wedge\tau})$ converges almost surely. By \eqref{eq:sum}, $\sum_{k<\tau}\gamma_kV(Y_k)<\infty$ almost surely, while $\sum_k\gamma_k=\infty$; so on $\{\tau=\infty\}$ the limit of $V(Y_k)$ is $0$. On the same event the series $\sum_k\gamma_k\mathcal Lq(Y_k)$ converges absolutely by \eqref{eq:box}, and $\mathcal M_K$ converges because it is bounded in $L^2$. Therefore $q_k\to X$ for some $X$ with $|X-\frac12|\le\eta$, and $\delta_k,r_k\to0$; that is, $u_k,v_k\to X$ and $z_k\to X(1-X)/2$. Replacing the denominator $n+2$ by $n$ does not change the limits, and the remaining normalised weights converge to $X$ and $1-X$ by the identities of Section~\ref{sec:drift}. For $X\in[49/100,51/100]$ we have $X(1-X)/2\in[2499/20000,1/8]$.

\subsection{Reaching the restart state}
Let $M$ be odd. Starting from unit weights, let the first $8M-2$ ants reinforce $U$ and $L$ each $(7M-1)/2$ times and $C$ and $B$ each $(M-1)/2$ times, in any fixed order. Every branch edge lies on two of the four paths with total count $(7M-1)/2+(M-1)/2=4M-1$, and $13$ lies on $C$ and $B$, so afterwards the branch weights are $4M$ and $W_{13}=M$. At every finite time each of the four paths has positive probability, since the walk can follow it step by step. So this history has positive probability, and for $M=50001$ the counts are $175003,175003,25000,25000$. By the Markov property and \eqref{eq:trap}, the event of Theorem~\ref{thm:main} has probability at least $3943/7143$ times the probability of this history, which is positive. This proves Theorem~\ref{thm:main}. \qed

\section{Other distances and the number of vertices}\label{sec:cor}

Corollary~\ref{cor:distance} and Remark~\ref{rem:minimal} follow from the block structure of spanning trees. Let $G$ be connected and let $\mathcal G_1,\dots,\mathcal G_s$ be the blocks of $G$ (the maximal $2$-connected subgraphs and the bridges) met by every path from $N$ to $F$, in order; let $N=c_0,c_1,\dots,c_s=F$ be the cut vertices at which consecutive blocks meet. A spanning tree of $G$ is the union of one spanning tree of each block, and its weight is the product of the weights of these trees. So under $\mathsf T$ the restrictions to the blocks are independent weighted spanning trees, and $\pi_{N,F}(\mathsf T)$ is the concatenation of the paths $\pi_{c_{i-1},c_i}$ in the trees of $\mathcal G_1,\dots,\mathcal G_s$. By Proposition~\ref{prop:same}, each ant therefore reinforces in each $\mathcal G_i$ a path whose conditional law is that of the ant process on $(\mathcal G_i,c_{i-1},c_i)$ with the current weights of $\mathcal G_i$. Hence the weights restricted to $\mathcal G_i$ form an ant process on $(\mathcal G_i,c_{i-1},c_i)$ started from unit weights, edges of other blocks are never reinforced, and a path from $N$ to $F$ is shortest exactly when each of its block pieces is shortest.

\begin{proof}[Proof of Corollary~\ref{cor:distance}]
Attach to $G_6$ a new path of $d-3$ edges ending at the vertex $0$, and put the nest at its other end. The new edges are bridges, so the weights on $G_6$ form the ant process of Theorem~\ref{thm:main}, while the distance from the nest to $5$ is $d$ and $13$ still lies on no shortest path.
\end{proof}

\begin{remark}\label{rem:minimal}
On a graph with at most five vertices the ants find the shortest paths. Indeed, a $2$-connected block with $b$ vertices has diameter at most $\lfloor b/2\rfloor$, since two vertices lie on a common cycle of length at most $b$; so every block met by the paths from $N$ to $F$ has $d(c_{i-1},c_i)\le2$. By the block structure, the argument of \cite[Proposition~2.8]{KMS22b} for distance one, \cite[Theorem~1.1]{AM26} for distance two and Proposition~\ref{prop:same}, every edge outside the shortest paths has normalised weight tending to zero under either rule. Together with Theorem~\ref{thm:main}, six is the least number of vertices of a counterexample.
\end{remark}

\section*{Acknowledgments and funding}

The research presented in this paper was supported by the European Research Council (ERC) under the European Union's Horizon Europe research and innovation programme (grant agreement No.~101041711) and by the Simons Foundation as part of the Collaboration on the Mathematical and Scientific Foundations of Deep Learning. It was also supported by Heights Labs, by Convex Nexus Capital and by the Israel Science Foundation (ISF grants 2258/19 and 4101/25).

This work was done by multiple AI agents, Claude (Anthropic) and Codex (OpenAI), connected to Math Brain~\cite{MathBrainRepo} and working together inside the Brain framework~\cite{BrainFramework}. The method is described in~\cite{MathBrainPaper}.

\section*{Declaration of competing interest}

The author holds equity in Heights Labs and in Convex Nexus Capital and has a management role at Convex Nexus Capital. These entities are included in the funding declaration above.

\section*{Declaration of generative AI and AI-assisted technologies}

This paper is part of a broader research project on proving open problems with AI systems. The results, proofs and text were produced with generative AI systems. OpenAI's Codex found the counterexample and wrote its proof, and a separate Codex run audited it. Anthropic's Claude (Opus~5.5) checked the proof and wrote the text. The models and the conjectures are due to Kious, Mailler and Schapira \cite{KMS22a} and to Armfield and Mailler \cite{AM26}. Proposition~\ref{prop:same} is a form of the time-reversal invariance of loop-erased random walk and rests on the spanning-tree algorithms of Wilson \cite{Wil96} and of Aldous \cite{Ald90} and Broder \cite{Bro89}. The author directed the work and is responsible for the content of the paper and for decisions concerning its dissemination.

\section*{Data and code availability}

The script \texttt{check.py}, available at \url{https://github.com/yspennstate/loop-erased-ants}, verifies in exact rational arithmetic the identity between the sum of the numerators in \eqref{eq:paths} and $D$, the drift \eqref{eq:drift}, the fourteen coefficients of \eqref{eq:DLV}, the constant $434717/9000000$ and the value $3943/7143$ in \eqref{eq:trap}. It also computes the law of the reinforced path on $G_6$ under both erasure rules at three weight vectors, the forward rule by the product formula for loop-erased random walk in terms of Green's functions and the backward rule by an exact absorbing Markov chain, and compares both with the spanning-tree law.

\begin{thebibliography}{99}

\bibitem{Ald90} D.~J.~Aldous, \emph{The random walk construction of uniform spanning trees and uniform labelled trees}, SIAM J. Discrete Math. \textbf{3} (1990), no.~4, 450--465.

\bibitem{AM26} W.~Armfield and C.~Mailler, \emph{The forward-loop-erased ants walk}, arXiv:2610.18453, 2026.

\bibitem{Bro89} A.~Z.~Broder, \emph{Generating random spanning trees}, in 30th Annual Symposium on Foundations of Computer Science, IEEE, 1989, 442--447.

\bibitem{KMS22a} D.~Kious, C.~Mailler and B.~Schapira, \emph{Finding geodesics on graphs using reinforcement learning}, Ann. Appl. Probab. \textbf{32} (2022), no.~5, 3889--3929.

\bibitem{KMS22b} D.~Kious, C.~Mailler and B.~Schapira, \emph{The trace-reinforced ants process does not find shortest paths}, J. \'Ec. polytech. Math. \textbf{9} (2022), 505--536. \url{https://doi.org/10.5802/jep.188}

\bibitem{LP16} R.~Lyons and Y.~Peres, \emph{Probability on Trees and Networks}, Cambridge Series in Statistical and Probabilistic Mathematics 42, Cambridge University Press, 2016.

\bibitem{MV26} C.~Mailler and Z.~Varin, \emph{The two-nest ants process on triangle-series-parallel graphs}, arXiv:2601.22855, 2026.

\bibitem{Wil96} D.~B.~Wilson, \emph{Generating random spanning trees more quickly than the cover time}, in Proceedings of the Twenty-Eighth Annual ACM Symposium on Theory of Computing, ACM, 1996, 296--303.

\bibitem{MathBrainRepo} Y.~Shmalo, \emph{Math Brain: public snapshot, code and tests}, 2026, \url{https://github.com/yspennstate/math-brain-public}.

\bibitem{MathBrainPaper} Y.~Shmalo, \emph{Math Brain: shared, auditable memory for multi-agent mathematical research}, October 2026, \url{https://github.com/yspennstate/math-brain-public/tree/main/papers/math-brain-2026-10-24}.

\bibitem{BrainFramework} Y.~Shmalo, \emph{Brain Framework}, public software repository, 2026, \url{https://github.com/yspennstate/brain-framework-public}.

\end{thebibliography}

\end{document}
